\begin{papercard}{gray2018}
\cf{Stages} \stagetag{C} \stagetag{W} (also \stagetag{H}). A practice-led content analysis of 118 practitioner-identified ``dark pattern'' artefacts. It first applies Brignull's typology deductively, then recodes openly, producing five \emph{strategy} categories: nagging, obstruction, sneaking, interface interference and forced action. \emph{Note:} PDF page 1 is a ResearchGate cover sheet.

\cf{Definition} Dark patterns are instances where designers use knowledge of human behaviour and users' desires to implement deceptive functionality that is not in the user's best interest \pc{gray2018}{2}{We use the term dark patterns to define instances where design}. Brignull's original 2010 definition is quoted, and the paper notes open questions about whose interests count and about intent \pc{gray2018}{4}{Which user or stakeholder interests are or should be kept in mind}. Brignull's typology mixes context, strategy and outcome, which makes comparison hard \pc{gray2018}{4}{the typology mixes context, strategy, and outcome, making comparison among patterns difficult}.

\cf{Corpus generation} Two researchers, one from UX and one from computer science (insider and outsider perspectives), collected iteratively for two months with weekly meetings to refine search strategies \pc{gray2018}{4}{selected one researcher from a user experience background}. Sources: search engines, social media (with a bot collecting all \#darkpatterns tweets), blogs, tech media and practitioner sites. The keyword set was broadened to ``unethical'', ``evil'' and ``manipulative'' combined with ``design'', ``UX'' and ``interface'' to find unlabelled cases \pc{gray2018}{5}{we identified key terms such as "un}. The 118 artefacts came from social media 45, blogs 40, news 19, own interactions 10 and darkpatterns.org 4. Each was catalogued with source, culprit site or app, a visual and a researcher memo \pc{gray2018}{5}{In total, we collected a corpus of 118 artifacts}.

\cf{Analysis} The constant comparative method and document analysis. Round 1: the Brignull categories as \emph{a priori} codes, which revealed overlaps and non-parallel categories \pc{gray2018}{5}{We first categorized artifacts using the Brignull [17] taxonomy as an a priori set of codes}. Round 2: open coding with emergent and axial codes, foregrounding interaction quality, context and user group. The result is a hierarchy organised by strategy and potential designer motivation, which subsumes Brignull's types \pc{gray2018}{5}{This led us to perform a second analysis of all exemplars in the corpus using an open coding approach}.

\cf{The five strategies (Figure 1)} \textbf{Nagging}: a minor redirection of expected function that persists across interactions, e.g.\ ``Not now''/``OK'' notification prompts with no permanent dismissal \pc{gray2018}{6}{We define nagging as a minor redirection of expected function}. \textbf{Obstruction}: a task flow made harder than needed to dissuade an action (Roach Motel, price-comparison prevention, intermediate currency), e.g.\ an ad-tracking opt-out hidden in iOS 6 settings with confusing wording \pc{gray2018}{6}{We define obstruction as impeding a task flow}. \textbf{Sneaking}: hiding, disguising or delaying relevant information (forced continuity, hidden costs, sneak into basket, bait and switch), e.g.\ consenting to data sale in order to unsubscribe \pc{gray2018}{7}{We define sneaking as an attempt to hide, disguise, or delay the divulging of information}. \textbf{Interface interference}: UI manipulation privileging some actions. Subtypes: \emph{hidden information}, \emph{preselection} \pc{gray2018}{8}{We define preselection as any situation where an option is selected by default}, and \emph{aesthetic manipulation} (toying with emotion, false hierarchy, disguised ads, trick questions) \pc{gray2018}{8}{False hierarchy gives one or more options visual or interactive precedence over others}. \textbf{Forced action}: a required action to access functionality (social pyramid, privacy zuckering, gamification) \pc{gray2018}{9}{We define forced action as any situation in which users are required to perform}.

\cf{Interpretation} Dark patterns undermine the user's interpretation of action possibilities (gulf of execution) and their expected outcomes (gulf of evaluation) \pc{gray2018}{9}{the user’s interpretation of the system and its communicative po}. Some dark patterns \emph{test well} in usability terms, at the expense of user choice \pc{gray2018}{9}{Paradoxically, some instances of dark patterns test well from a usability perspective}. Several parallel Fogg's persuasive strategies, e.g.\ tunnelling and reduction parallel forced action and obstruction.

\cf{Use for this thesis} The five strategy codes plus subtypes form a ready-made \emph{deductive code family} for walkthrough observations in consent, settings and deletion flows: preselection in onboarding, obstruction of deletion, nagging for permissions, hidden information in long terms. The method is itself a model of hybrid coding: start with an \emph{a priori} typology, document where it fails, then recode openly and restructure. This is useful for justifying codebook revisions in the pilot. Caution: the corpus is a convenience sample of practitioner-flagged cases, so it gives category definitions, not prevalence. ``Designer intent'' cannot be observed in a walkthrough, so code the \emph{interface effect} and avoid attributing intent.
\end{papercard}
